\chapter{Assembly and the Build Pipeline}
\label{chap:build}

A finished operating system is not a single program. It is a carefully arranged
constellation of artifacts --- a static C \code{init} that the kernel hands
control to as PID~1, a Go multicall binary that pretends to be a dozen
familiar coreutils, a Rust package manager, a sprinkling of kernel modules,
configuration files in their canonical FHS locations, a compressed root
filesystem image, and a bootloader that knows how to find all of it. The
interesting engineering question is not how to compile any one of these
pieces, but how to make a single command turn a polyglot source tree into a
bootable image, reproducibly and in the right order.

This chapter walks the \emph{actual} AULinux build pipeline as it exists in
the \file{scripts/} directory. Every command, flag, and path quoted below is
taken verbatim from the scripts; where a script is partial, assumes host
prerequisites, or relies on mock content, that is called out honestly rather
than papered over.

\section{The Challenge of a Polyglot Build}

AULinux is written in four languages, each with its own toolchain and its own
opinion about output layout:

\begin{itemize}
  \item The \textbf{init system} is C, compiled with \code{gcc} through a
        hand-written \file{Makefile}.
  \item The \textbf{userland} (shell plus coreutils) is Go, compiled into a
        single \emph{multicall} binary with \code{go build}.
  \item The \textbf{package manager} is Rust, compiled with
        \code{cargo build --release}.
  \item The \textbf{kernel modules} are C, compiled out-of-tree against the
        host's kbuild infrastructure.
\end{itemize}

No single language's build tool can orchestrate the others. Cargo does not
know how to invoke \code{go build}; \code{make} does not understand Cargo
target triples; none of them know how to lay out an FHS tree or pack a cpio
archive. The pragmatic glue layer is therefore a set of POSIX shell scripts.
Shell is the lowest common denominator that can call every toolchain, copy
files, create symlinks, and shell out to \code{cpio}, \code{tar}, and
\code{qemu} --- exactly the heterogeneous, side-effect-heavy work that
operating-system assembly consists of.

\begin{infobox}[The orchestration model]
The pipeline is staged on purpose. \file{build.sh} \emph{compiles} each
subsystem into \file{build/}, knowing nothing about filesystems or boot.
\file{make\_rootfs.sh} then \emph{assembles} those artifacts into an FHS tree,
and \file{run\_qemu.sh} \emph{packs and boots} it. Compilation, assembly, and
boot are cleanly separated, so any stage can be re-run independently.
\end{infobox}

\begin{table}[h]
\centering
\caption{Each subsystem, its toolchain, build command, and output artifact.}
\begin{tabular}{@{}llll@{}}
\toprule
\textbf{Subsystem} & \textbf{Language} & \textbf{Build command} & \textbf{Artifact} \\
\midrule
Init       & C    & \code{make} (gcc \code{-static -O2})                        & \file{build/init/au-init} \\
Utils + Shell & Go & \code{go build} (\code{CGO\_ENABLED=0}, \code{-s -w -extldflags '-static'}) & \file{build/utils/auutils} (+ symlinks) \\
Package mgr & Rust & \code{cargo build --release} (\code{crt-static}, x86\_64 target) & \file{build/pkg-manager/aupkg} \\
Kernel mods & C   & \code{make -C \$KDIR M=src modules}                         & \file{build/kernel/*.ko} \\
\bottomrule
\end{tabular}
\end{table}

\section{\file{build.sh}: The Orchestrator}

\file{build.sh} is the entry point. It begins defensively --- \code{set -e} on
line~7 makes any non-zero exit abort the whole build --- and resolves the
project root from its own location, so it works regardless of the caller's
working directory:

\begin{lstlisting}[language=bash, caption={build.sh --- safe root resolution and fail-fast mode}]
set -e
# Project root
PROJECT_ROOT="$(cd "$(dirname "$0")/.." && pwd)"
BUILD_DIR="$PROJECT_ROOT/build"
\end{lstlisting}

The \code{all} target (the default, dispatched by the \code{case} at the
bottom of the file) calls the per-subsystem builders in a deliberate order:

\begin{lstlisting}[language=bash, caption={build.sh --- the build\_all driver and its phase order}]
build_all() {
    print_banner
    mkdir -p "$BUILD_DIR"

    build_init
    # build_shell is covered by build_utils now
    build_utils
    build_pkgmanager
    build_kernel
    ...
}
\end{lstlisting}

The real order is therefore: \textbf{init $\rightarrow$ utils (which is also
the shell) $\rightarrow$ package manager $\rightarrow$ kernel modules}. Note
that the shell is no longer built separately; \code{build\_shell} simply
delegates to \code{build\_utils}, because the shell \code{aush} is now one of
the applets inside the Go multicall binary.

\subsection{Phase 1 --- Init (C, statically linked)}

\code{build\_init} is the simplest phase: it changes into \file{init/} and runs
\code{make}. The \file{init/Makefile} compiles \file{src/init.c} with
\code{gcc} and the flags \code{-Wall -Wextra -O2 -static}, emitting
\file{build/init/au-init}.

\begin{infobox}[Why \code{-static} for init]
PID~1 is the very first userspace program the kernel executes from the
initramfs. At that instant no dynamic loader has been configured and no shared
libraries are guaranteed to be present. A statically linked \code{au-init}
carries its entire libc inside the executable, so it is fully relocatable and
cannot fail with a ``library not found'' error during early boot. This is the
single most important linking decision in the whole image.
\end{infobox}

\subsection{Phase 2 --- Utilities and shell (Go multicall binary)}

This is the most carefully optimized phase. A multicall binary is a single
executable that inspects \code{argv[0]} (the name it was invoked as) and
dispatches to the matching applet --- the same trick BusyBox uses. One binary
provides \code{ls}, \code{cat}, \code{cp}, the \code{aush} shell, the
\code{aubuild} package builder, and more, with the per-applet cost being a
single symlink rather than a whole executable.

\begin{lstlisting}[language=bash, caption={build.sh --- compiling and shrinking the auutils multicall binary}]
# Build from root to pick up go.mod
cd "$PROJECT_ROOT"

info "Compiling auutils..."
CGO_ENABLED=0 go build -ldflags="-s -w -buildid= -extldflags '-static'" \
    -trimpath -o "$BUILD_DIR/utils/auutils" utils/main.go

# Compression (Optional)
if command -v upx >/dev/null; then
    info "Compressing binary with UPX..."
    upx --best --lzma "$BUILD_DIR/utils/auutils"
else
    warn "UPX not found, skipping compression."
fi
\end{lstlisting}

Every flag here is earning its place:

\begin{itemize}
  \item \code{CGO\_ENABLED=0} disables cgo entirely, so the Go toolchain never
        calls out to the host C compiler and produces a pure-Go object graph with
        no implicit dependency on the host's \code{libc.so}.
  \item \code{-ldflags="-s -w"} strips the symbol table (\code{-s}) and DWARF
        debug information (\code{-w}), shaving megabytes off the binary.
  \item \code{-buildid=} blanks the build ID, reducing build-to-build variance.
  \item \code{-extldflags '-static'} passes \code{-static} to the external
        linker, forcing fully static linking at the linker level. Combined with
        \code{CGO\_ENABLED=0}, this guarantees the resulting binary has zero
        shared-library dependencies --- it will run on any \code{x86\_64} Linux
        system without a dynamic loader or libc present.
  \item \code{-trimpath} removes absolute filesystem paths from the binary,
        which both reduces size and improves reproducibility.
  \item \code{upx --best --lzma} (when present) further compresses the
        executable, decompressing transparently at launch.
\end{itemize}

After compilation, the script builds a \emph{link farm} --- one relative
symlink per applet, all pointing at \code{auutils}:

\begin{lstlisting}[language=bash, caption={build.sh --- creating the applet symlink farm}]
local UTILS=("ls" "cat" "cp" "mv" "mkdir" "rm" "chmod" "echo" "pwd" "ps" \
             "ip" "useradd" "groupadd" "aubuild" "aush" "sh")
info "Creating symlinks..."
for tool in "${UTILS[@]}"; do
    ln -sf auutils "$BUILD_DIR/utils/$tool"
done
\end{lstlisting}

When the rootfs later invokes \file{/bin/ls}, the kernel follows the symlink to
\code{auutils}, which reads \code{argv[0]}, sees \code{ls}, and runs the
listing applet. Note that \code{sh} is symlinked here too, so the system shell
is the same binary.

\subsection{Phase 3 --- Package manager (Rust, fully static)}

The Rust phase builds an optimized, statically-linked \code{aupkg} for an
explicit target triple and then copies it out of Cargo's deep target tree into
the flat \file{build/} layout the rest of the pipeline expects:

\begin{lstlisting}[language=bash, caption={build.sh --- static release build of the aupkg package manager}]
build_pkgmanager() {
    info "Building package manager (aupkg)..."
    cd "$PROJECT_ROOT/pkg-manager"
    RUSTFLAGS="-C target-feature=+crt-static" \
        cargo build --release --target x86_64-unknown-linux-gnu
    mkdir -p "$BUILD_DIR/pkg-manager"
    cp target/x86_64-unknown-linux-gnu/release/aupkg "$BUILD_DIR/pkg-manager/"
    success "Package manager built"
}
\end{lstlisting}

Two additions beyond a plain \code{cargo build --release} make the binary
fully self-contained:

\begin{itemize}
  \item \code{RUSTFLAGS="-C target-feature=+crt-static"} instructs the Rust
        compiler to link the C runtime (\code{crt}) statically rather than
        depending on a dynamic \code{libc.so} at runtime. The result is a binary
        that carries its own C runtime, with no dependency on the host's glibc
        version.
  \item \code{--target x86\_64-unknown-linux-gnu} pins the Cargo target triple
        explicitly. This means Cargo places the artifact in
        \file{target/x86\_64-unknown-linux-gnu/release/aupkg} rather than the
        default \file{target/release/aupkg}, and the copy command reflects that
        deeper path. An explicit triple also ensures the build is reproducible
        regardless of the host machine's default Rust target.
\end{itemize}

Together these choices give \code{aupkg} the same ``runs anywhere on
\code{x86\_64}'' property as \code{au-init} and \code{auutils}.

\subsection{Phase 4 --- Kernel modules (out-of-tree C)}

\code{build\_kernel} is the only phase allowed to fail without aborting the
build. The \file{kernel/Makefile} uses kbuild to compile three modules ---
\code{aulinux\_core}, \code{aulinux\_sysctl}, and \code{aulinux\_procmon} ---
against \code{KDIR} (default \code{/lib/modules/\$(uname -r)/build}). Because
that requires the host's kernel headers, the orchestrator wraps it
defensively:

\begin{lstlisting}[language=bash, caption={build.sh --- tolerant kernel-module phase}]
build_kernel() {
    info "Building kernel modules..."
    cd "$PROJECT_ROOT/kernel"
    if make; then
        success "Kernel modules built"
    else
        warn "Kernel modules build failed (may require kernel headers)"
    fi
}
\end{lstlisting}

\begin{notebox}[Kernel headers are a host prerequisite]
The module build only succeeds on a host that has matching kernel headers
installed at \code{/lib/modules/\$(uname -r)/build}. On a headers-less machine
the phase emits a warning and the build continues --- the resulting image
simply ships without the \code{.ko} files. The modules are \emph{not} required
to boot the QEMU image, which uses a separately obtained kernel (see below).
\end{notebox}

\section{Constructing the Root Filesystem (\file{make\_rootfs.sh})}

With \file{build/} populated, \file{make\_rootfs.sh} assembles a Filesystem
Hierarchy Standard (FHS) tree under \file{rootfs/}. It too runs under
\code{set -e}.

\subsection{The FHS skeleton}

It first creates the conventional top-level directories --- now including
\file{lib} and \file{lib64} which are needed by the dynamic loader and shared
library machinery --- and tightens permissions on \file{/tmp} (the sticky bit,
\code{1777}):

\begin{lstlisting}[language=bash, caption={make\_rootfs.sh --- the FHS directory skeleton}]
mkdir -p "$ROOTFS_DIR"/{bin,sbin,etc,usr,var,proc,sys,dev,tmp,run,boot,home,root,lib,lib64}
mkdir -p "$ROOTFS_DIR"/etc/aulinux
mkdir -p "$ROOTFS_DIR"/usr/{bin,sbin,lib}
chmod 1777 "$ROOTFS_DIR/tmp"
\end{lstlisting}

\subsection{Shared libraries and the dynamic loader}

Even though the AULinux-authored binaries (\code{au-init}, \code{auutils},
\code{aupkg}) are all fully statically linked, the rootfs also ships several
\emph{host} utilities --- \code{tar}, \code{gzip}, \code{dhcpcd} --- that are
dynamically linked against glibc. These binaries require the dynamic loader
\code{ld-linux-x86-64.so.2} and a set of \code{.so} files to be present inside
the rootfs itself, because there is no host filesystem to fall back on once the
initramfs is running.

The script copies the loader and the required libraries using \code{cp -L},
which dereferences symlinks so the real versioned \code{.so} files land in
\file{rootfs/lib/} rather than dangling symlinks:

\begin{lstlisting}[language=bash, caption={make\_rootfs.sh --- loader, core glibc libs, and host-utility libs}]
cp -L /lib64/ld-linux-x86-64.so.2       "$ROOTFS_DIR/lib64/"
cp -L /lib/x86_64-linux-gnu/libc.so.6   "$ROOTFS_DIR/lib/"
cp -L /lib/x86_64-linux-gnu/libm.so.6   "$ROOTFS_DIR/lib/"
cp -L /lib/x86_64-linux-gnu/libresolv.so.2  "$ROOTFS_DIR/lib/"
cp -L /lib/x86_64-linux-gnu/librt.so.1  "$ROOTFS_DIR/lib/"
cp -L /lib/x86_64-linux-gnu/libdl.so.2  "$ROOTFS_DIR/lib/"
cp -L /lib/x86_64-linux-gnu/libpthread.so.0 "$ROOTFS_DIR/lib/"

# Extra libraries required by host utilities (like tar)
cp -L /usr/lib/x86_64-linux-gnu/libacl.so.1    "$ROOTFS_DIR/lib/"
cp -L /usr/lib/x86_64-linux-gnu/libselinux.so.1 "$ROOTFS_DIR/lib/"
cp -L /usr/lib/x86_64-linux-gnu/libpcre2-8.so.0 "$ROOTFS_DIR/lib/"

# Canonical compatibility symlink so /lib and /lib64 agree on the loader
ln -sf /lib64/ld-linux-x86-64.so.2 "$ROOTFS_DIR/lib/ld-linux-x86-64.so.2"
\end{lstlisting}

The following table lists every library copied in this section and why it is
needed:

\begin{table}[h]
\centering
\caption{Shared libraries copied into \file{rootfs/lib/} and their purpose.}
\begin{tabular}{@{}ll@{}}
\toprule
\textbf{Library} & \textbf{Purpose} \\
\midrule
\file{ld-linux-x86-64.so.2} & ELF dynamic linker/loader --- required by every dynamically linked binary \\
\file{libc.so.6}             & GNU C library --- core POSIX and C standard library \\
\file{libm.so.6}             & Math library (\code{sin}, \code{cos}, \code{sqrt}, \ldots) \\
\file{libresolv.so.2}        & DNS resolver routines used by networking code \\
\file{librt.so.1}            & POSIX real-time extensions (clocks, message queues) \\
\file{libdl.so.2}            & \code{dlopen}/\code{dlsym} --- runtime shared-library loading \\
\file{libpthread.so.0}       & POSIX threads \\
\file{libacl.so.1}           & POSIX ACL support, required by \code{tar} \\
\file{libselinux.so.1}       & SELinux label APIs, linked by \code{tar} on SELinux-capable hosts \\
\file{libpcre2-8.so.0}       & PCRE2 regex library, a dependency of \code{libselinux} \\
\bottomrule
\end{tabular}
\end{table}

\begin{warnbox}[Host-architecture coupling]
These paths are hard-coded for a Debian/Ubuntu \code{x86\_64} host
(\file{/lib/x86\_64-linux-gnu/...}). On a different distribution, a non-glibc
host, or a non-\code{x86\_64} architecture, these copies will fail under
\code{set -e} and abort \file{make\_rootfs.sh}. The rootfs is not portable
across host layouts as written.
\end{warnbox}

\subsection{Host archivers: \code{tar} and \code{gzip}}

After the shared libraries are in place, the script copies two host archiving
utilities directly into \file{rootfs/bin/}:

\begin{lstlisting}[language=bash, caption={make\_rootfs.sh --- copying host tar and gzip into the rootfs}]
cp -L /usr/bin/tar  "$ROOTFS_DIR/bin/"
cp -L /usr/bin/gzip "$ROOTFS_DIR/bin/"
\end{lstlisting}

\code{aupkg} relies on \code{tar} and \code{gzip} to unpack \file{.tar.gz}
package archives at install time. Rather than building these tools from source,
the script reuses the host binaries, which are dynamically linked --- hence the
glibc and \code{libacl}/\code{libselinux}/\code{libpcre2} libraries copied
above are a prerequisite for \code{tar} to function inside the initramfs.

\subsection{DHCP client setup}

Networking support is delivered by copying a real DHCP client from the host
system. The script probes four possible locations in order:

\begin{lstlisting}[language=bash, caption={make\_rootfs.sh --- probing for and copying the DHCP client}]
if [ -f "/usr/sbin/dhcpcd" ]; then   DHCP_CLIENT="/usr/sbin/dhcpcd"
elif [ -f "/sbin/dhcpcd" ]; then     DHCP_CLIENT="/sbin/dhcpcd"
elif [ -f "/sbin/dhclient" ]; then   DHCP_CLIENT="/sbin/dhclient"
elif [ -f "/usr/sbin/dhclient" ]; then DHCP_CLIENT="/usr/sbin/dhclient"
fi

if [ -n "$DHCP_CLIENT" ]; then
    cp -L "$DHCP_CLIENT" "$ROOTFS_DIR/sbin/dhcpcd"
    chmod +x "$ROOTFS_DIR/sbin/dhcpcd"
    # Crypto and compression libraries that dhcpcd links against
    [ -f "/usr/lib/x86_64-linux-gnu/libcrypto.so.3" ] && \
        cp -L /usr/lib/x86_64-linux-gnu/libcrypto.so.3 "$ROOTFS_DIR/lib/"
    [ -f "/usr/lib/x86_64-linux-gnu/libz.so.1" ] && \
        cp -L /usr/lib/x86_64-linux-gnu/libz.so.1 "$ROOTFS_DIR/lib/"
    [ -f "/usr/lib/x86_64-linux-gnu/libzstd.so.1" ] && \
        cp -L /usr/lib/x86_64-linux-gnu/libzstd.so.1 "$ROOTFS_DIR/lib/"
fi
\end{lstlisting}

Whichever client is found, it is installed as \file{/sbin/dhcpcd} inside the
rootfs, normalising the name so init's service table can refer to it by a
single well-known path. The additional libraries (\code{libcrypto},
\code{libz}, \code{libzstd}) are guarded with \code{[ -f ... ]} so a missing
library does not abort the build --- they are only needed if \code{dhcpcd}
actually links against them on the particular host.

\subsection{The \code{services.conf} service table and DNS}

AULinux's init reads \file{/etc/aulinux/services.conf} to learn which
long-running daemons to supervise. Previously this file was written as an
empty stub. It now contains a real entry:

\begin{lstlisting}[language=bash, caption={make\_rootfs.sh --- writing services.conf with the DHCP client entry}]
cat > "$ROOTFS_DIR/etc/aulinux/services.conf" << EOF
# AULinux Services Configuration
# Format: name=executable args
dhcpcd=/sbin/dhcpcd -B
EOF
\end{lstlisting}

The \code{-B} flag runs \code{dhcpcd} in the \emph{foreground} (no fork-to-
background), which is the correct mode for a process supervised by an init
system --- init can observe the process, detect exits, and restart it.

DNS resolution is configured by writing a static \file{/etc/resolv.conf} with
two public resolvers:

\begin{lstlisting}[language=bash, caption={make\_rootfs.sh --- configuring public DNS via resolv.conf}]
cat > "$ROOTFS_DIR/etc/resolv.conf" << EOF
nameserver 8.8.8.8
nameserver 1.1.1.1
EOF
\end{lstlisting}

Together, the DHCP client entry in \file{services.conf} (which acquires a
dynamic IP address on \code{eth0}) and the static \file{resolv.conf} (which
provides name resolution) give AULinux a functional TCP/IP stack on first boot
with no manual configuration.

\subsection{Wiring \code{/init}}

The kernel, after unpacking the initramfs, executes \file{/init} by default.
AULinux installs the C init binary to \file{/sbin/init} and then creates a
top-level \file{/init} symlink pointing at it, satisfying both the initramfs
convention and the more traditional \file{/sbin/init} location:

\begin{lstlisting}[language=bash, caption={make\_rootfs.sh --- installing init and the /init entry point}]
if [ -f "$BUILD_DIR/init/au-init" ]; then
    cp "$BUILD_DIR/init/au-init" "$ROOTFS_DIR/sbin/init"
    ln -sf /sbin/init "$ROOTFS_DIR/init"
    echo "Installed init"
else
    echo "Warning: au-init not found in build/"
fi
\end{lstlisting}

\subsection{Placing userland, shell, and package manager}

The whole \file{build/utils/} directory --- the \code{auutils} binary
\emph{and} its symlink farm --- is copied into \file{/bin} with \code{cp -d}
so the symlinks are preserved rather than dereferenced. The package manager
goes to \file{/usr/bin/aupkg} with a convenience symlink in \file{/bin}. The
shell needs no separate copy: it already arrived as the \code{aush}/\code{sh}
symlinks inside \file{build/utils/}, and the script explicitly recognizes this
case (``Shell is integrated into utilities'').

\begin{lstlisting}[language=bash, caption={make\_rootfs.sh --- installing utils (symlink-preserving) and aupkg}]
if [ -d "$BUILD_DIR/utils" ]; then
    cp -d "$BUILD_DIR/utils/"* "$ROOTFS_DIR/bin/" 2>/dev/null || true
    echo "Installed utilities"
fi

if [ -f "$BUILD_DIR/pkg-manager/aupkg" ]; then
    cp "$BUILD_DIR/pkg-manager/aupkg" "$ROOTFS_DIR/usr/bin/aupkg"
    ln -sf /usr/bin/aupkg "$ROOTFS_DIR/bin/aupkg"
    echo "Installed package manager"
fi
\end{lstlisting}

\subsection{System configuration}

The script then writes the minimal set of configuration files a Unix system
expects --- \file{/etc/hostname}, \file{/etc/passwd}, \file{/etc/group}, and
\file{/etc/fstab} --- guarding the account files with \code{[ ! -f ... ]} so a
hand-customized version is never clobbered:

\begin{lstlisting}[language=bash, caption={make\_rootfs.sh --- baseline accounts and mount table}]
echo "aulinux" > "$ROOTFS_DIR/etc/hostname"

cat > "$ROOTFS_DIR/etc/passwd" << EOF
root:x:0:0:root:/root:/bin/aush
guest:x:1000:1000:guest:/home/guest:/bin/aush
EOF

cat > "$ROOTFS_DIR/etc/fstab" << EOF
proc             /proc         proc    defaults   0    0
sysfs            /sys          sysfs   defaults   0    0
EOF
\end{lstlisting}

Both accounts use \file{/bin/aush} as their login shell --- the same multicall
binary. The \file{fstab} declares the two pseudo-filesystems init will mount.

\subsection{Seeding the package repository and the desktop/installer layer}

Finally, \file{make\_rootfs.sh} populates an on-image package repository under
\file{/var/lib/aupkg} and installs optional desktop and installer components.
Each invocation is guarded by a \code{[ -f ... ]} test so the build continues
even if a helper script is absent:

\begin{lstlisting}[language=bash, caption={make\_rootfs.sh --- package repo population and optional desktop/installer layer}]
# 1. Mock graphics skeleton
if [ -f "$PROJECT_ROOT/scripts/populate_graphics_repo.sh" ]; then
    "$PROJECT_ROOT/scripts/populate_graphics_repo.sh"
fi
# 2. Overwrite with real host-extracted graphics binaries
if [ -f "$PROJECT_ROOT/scripts/build_real_graphics_packages.py" ]; then
    python3 "$PROJECT_ROOT/scripts/build_real_graphics_packages.py"
fi
# 3. Networking and SSH packages
if [ -f "$PROJECT_ROOT/scripts/build_networking_packages.py" ]; then
    python3 "$PROJECT_ROOT/scripts/build_networking_packages.py"
fi
# 4. Desktop environments (GNOME, Plasma, ...)
if [ -f "$PROJECT_ROOT/scripts/build_desktop_environments.py" ]; then
    python3 "$PROJECT_ROOT/scripts/build_desktop_environments.py"
fi
# 5. Desktop launcher and OS installer wizard
if [ -f "$PROJECT_ROOT/scripts/start_desktop.sh" ]; then
    cp "$PROJECT_ROOT/scripts/start_desktop.sh" "$ROOTFS_DIR/bin/start_desktop.sh"
    chmod +x "$ROOTFS_DIR/bin/start_desktop.sh"
    ln -sf start_desktop.sh "$ROOTFS_DIR/bin/start_desktop"
fi
if [ -f "$PROJECT_ROOT/scripts/install_os.sh" ]; then
    cp "$PROJECT_ROOT/scripts/install_os.sh" "$ROOTFS_DIR/bin/au-install"
    chmod +x "$ROOTFS_DIR/bin/au-install"
    ln -sf au-install "$ROOTFS_DIR/bin/install-os"
fi
\end{lstlisting}

\begin{notebox}[The seeded packages are mock content]
\file{populate\_graphics\_repo.sh} and the \file{build\_*.py} helpers do not
compile real upstream software. They construct \emph{mock} package payloads
--- small shell stubs and placeholder \code{.so} files --- \code{tar.gz} them
into \file{/var/lib/aupkg/packages/}, and merge rich JSON metadata
(\code{name}, \code{version}, \code{dependencies}, \code{license},
\code{installed\_size}, \ldots) into \file{repo.db}. Their value is exercising
\code{aupkg}'s dependency resolution and install machinery end-to-end (for
example \code{openbox} $\rightarrow$ \code{xorg-server} $\rightarrow$
\code{libinput}), not shipping a working X11 stack. The repository covers a
graphics layer (libinput, xorg-server, xf86-video-fbdev, openbox), GUI
libraries and apps (cairo, pango, libgtk-3, firefox, vscode), a Wayland stack
(libwayland, wayland-protocols, weston), an XFCE desktop (xfce4-session,
thunar, alacritty), desktop environments (GNOME, Plasma), and networking tools
(busybox, dropbear, openssh).
\end{notebox}

\section{Booting in QEMU (\file{run\_qemu.sh})}

\file{run\_qemu.sh} is the one-command path from source tree to a running
system. It is fully automated and chains every prior stage. Its phases are:

\begin{enumerate}
  \item \textbf{Build} --- runs \code{scripts/build.sh all}.
  \item \textbf{Rootfs} --- deletes and regenerates \file{rootfs/} via
        \code{make\_rootfs.sh}.
  \item \textbf{Kernel acquisition} --- AULinux ships no kernel of its own, so
        the script tries, in order: a pre-downloaded Alpine
        \file{vmlinuz-virt}; the host's \file{/boot/vmlinuz}; the versioned
        \file{/boot/vmlinuz-\$(uname -r)}; an \code{apt-get download} of the
        host kernel package (extracted with \code{dpkg-deb -x}); and finally a
        download of a minimal Alpine virtual kernel over the network.
  \item \textbf{Initramfs packaging} --- packs \file{rootfs/} into a gzipped
        cpio archive.
  \item \textbf{Boot} --- launches QEMU.
\end{enumerate}

\subsection{Why initramfs, and how it is packed}

The boot image is an \textbf{initramfs}, not an ext4 disk. An initramfs is a
\code{cpio} archive that the kernel unpacks into an in-memory \code{tmpfs} and
then executes \file{/init} from directly --- there is no block device, no
mounting, no filesystem driver required to reach userspace. That makes it the
simplest possible thing to boot under an emulator. The archive uses the
\code{newc} (SVR4) cpio format the kernel expects, gzip-compressed:

\begin{lstlisting}[language=bash, caption={run\_qemu.sh --- packing the rootfs into a gzipped newc cpio initramfs}]
info "Packaging rootfs into initramfs.img..."
cd "$ROOTFS_DIR"
find . -print0 | cpio --null -ov --format=newc | gzip -9 > "$INITRAMFS_FILE"
cd "$PROJECT_ROOT"
\end{lstlisting}

\begin{infobox}[initramfs vs. rootfs-on-disk]
The same \file{rootfs/} tree feeds two different boot strategies. Packed as a
cpio initramfs (this script) it boots entirely from RAM --- ideal for QEMU.
Copied onto an ext4 partition (\file{install\_grub.sh}, next section) it
becomes a persistent on-disk root filesystem. AULinux supports both from one
source of truth.
\end{infobox}

\subsection{The QEMU invocation}

QEMU is given the kernel and the initramfs directly, bypassing any bootloader.
There are two modes selected by a \code{--gui} flag:

\begin{lstlisting}[language=bash, caption={run\_qemu.sh --- serial-console (default) and graphical QEMU launches}]
# Default: serial console, no window
qemu-system-x86_64 \
    -kernel "$KERNEL_FILE" \
    -initrd "$INITRAMFS_FILE" \
    -append "init=/sbin/init console=ttyS0 quiet" \
    -nographic \
    -m 512M

# With --gui: graphical output
qemu-system-x86_64 \
    -kernel "$KERNEL_FILE" \
    -initrd "$INITRAMFS_FILE" \
    -append "init=/sbin/init quiet" \
    -m 512M
\end{lstlisting}

Each flag does specific work:

\begin{itemize}
  \item \code{-kernel} loads the bzImage directly (QEMU acts as a Linux
        loader, so no GRUB is needed for emulation).
  \item \code{-initrd} hands the kernel the gzipped cpio as the initial
        ramdisk.
  \item \code{-append} sets the kernel command line. \code{init=/sbin/init}
        tells the kernel which program to run as PID~1 (overriding the
        \file{/init} default); \code{console=ttyS0} redirects the console to
        the first serial port so it appears in the terminal; \code{quiet}
        suppresses verbose kernel logging.
  \item \code{-nographic} (default mode) drops the graphical window and wires
        the VM's serial port to the host terminal --- the script reminds the
        user to exit with \code{Ctrl+A} then \code{X}.
  \item \code{-m 512M} allocates 512\,MiB of guest RAM, which must comfortably
        hold the decompressed initramfs plus shared libraries and archivers.
\end{itemize}

\subsection{Assembling a desktop session}

When a graphical session is wanted, \file{start\_desktop.sh} (installed into
\file{/bin}) orchestrates it from inside the running system. It exports
\code{PATH}, \code{LD\_LIBRARY\_PATH}, \code{HOME}, and \code{DISPLAY=:0}; uses
\code{aupkg install openbox xf86-video-fbdev} to pull in the graphics stack
(relying on \code{aupkg}'s dependency resolution to also bring in
\code{xorg-server} and \code{libinput}); starts \code{Xorg :0} against an
\code{fbdev} configuration; and finally launches the \code{openbox} window
manager. Because the seeded packages are mock stubs (see the note above), this
demonstrates the install-and-launch \emph{flow} rather than rendering a real
desktop.

\section{Bootable Live ISO (\file{build\_iso.sh})}

\file{build\_iso.sh} extends the QEMU path into a distributable, bootable ISO
image suitable for writing to USB media or running in a full VM. It requires
\code{grub-mkrescue} (from \code{grub-common} / \code{grub-pc-bin}) and
\code{xorriso}, and chains every prior stage automatically:

\begin{enumerate}
  \item Calls \code{build.sh all} to compile all components.
  \item Calls \code{make\_rootfs.sh} to assemble a fresh \file{rootfs/}.
  \item Obtains a kernel using the same waterfall as \file{run\_qemu.sh}
        (pre-downloaded, host \file{/boot/vmlinuz}, versioned fallback,
        or Alpine download).
  \item Packs \file{rootfs/} into \file{build/initramfs.img} via
        \code{find | cpio | gzip}.
  \item Stages \file{boot/vmlinuz} and \file{boot/initramfs.img} under a
        temporary \file{build/iso\_tmp/} tree alongside a custom GRUB
        configuration.
  \item Calls \code{grub-mkrescue -o build/aulinux.iso build/iso\_tmp/}.
\end{enumerate}

The embedded \file{grub.cfg} offers three boot menu entries:

\begin{lstlisting}[language=bash, caption={build\_iso.sh --- GRUB menu entries embedded in the ISO}]
menuentry "AULinux Live OS (Console Mode)" {
    linux /boot/vmlinuz init=/sbin/init console=ttyS0 console=tty0 quiet
    initrd /boot/initramfs.img
}
menuentry "AULinux Live OS (Graphical Mode - Openbox)" {
    linux /boot/vmlinuz init=/sbin/init quiet
    initrd /boot/initramfs.img
}
menuentry "Install AULinux to Hard Disk (Interactive Wizard)" {
    linux /boot/vmlinuz init=/bin/install-os quiet
    initrd /boot/initramfs.img
}
\end{lstlisting}

The third entry boots directly into \file{/bin/install-os} (the
\code{install\_os.sh} wizard) as PID~1 instead of the normal init, providing a
dedicated installer path on the live media.

\section{The OS Installation Wizard (\file{install\_os.sh})}

\file{install\_os.sh} is the interactive installation wizard that the
\code{build\_iso.sh} live media boots into when the user selects ``Install
AULinux to Hard Disk''. It is copied into \file{rootfs/bin/au-install} (with a
\file{/bin/install-os} symlink) by \file{make\_rootfs.sh}. The wizard walks
through the canonical installation steps:

\begin{itemize}
  \item Scans and lists available storage drives.
  \item Accepts a target disk selection (defaulting to \code{/dev/sda}).
  \item Partitions the target with a single bootable ext4 partition.
  \item Formats the partition and reports its UUID.
  \item Copies the AULinux filesystem contents to \file{/mnt} in stages
        (dynamic loader and glibc, utilities, init, \code{aupkg} and its
        repository database, package archives).
  \item Installs GRUB to the target MBR and writes \file{/mnt/boot/grub/grub.cfg}.
\end{itemize}

\begin{notebox}[The wizard is a demonstrator, not a production installer]
\file{install\_os.sh} is written in plain \code{/bin/sh} and simulates each
step with \code{sleep} delays and \code{echo} output rather than calling
\code{parted}, \code{mkfs.ext4}, or \code{grub-install} for real. Its purpose
is to exercise the full boot path --- booting the ISO, reaching the installer
as PID~1, and displaying a credible installation flow --- rather than to write
a working disk. A real installer would replace the \code{echo}/\code{sleep}
stubs with the actual partitioning commands.
\end{notebox}

\section{Bootloader Installation (\file{install\_grub.sh})}

For a persistent, bootable-on-real-hardware (or full-VM) image, GRUB must be
installed onto a disk image. \file{install\_grub.sh} targets a 512\,MiB
\code{msdos}/ext4 image. Crucially, it is \textbf{semi-manual by design}: it
prints the exact sequence of root-only, loop-device operations and only
executes them automatically if it detects it is running as root.

\begin{lstlisting}[language=bash, caption={install\_grub.sh --- the documented disk-image and GRUB sequence}]
# 1. Create disk image:
dd if=/dev/zero of=$IMAGE_FILE bs=1M count=512
# 2. Partition image:
parted $IMAGE_FILE --script mklabel msdos mkpart primary ext4 1MiB 100% set 1 boot on
# 3. Format & mount via loop device:
LOOPDEV=$(losetup -fP --show $IMAGE_FILE)
mkfs.ext4 ${LOOPDEV}p1 ; mount ${LOOPDEV}p1 /mnt
# 4. Copy rootfs:    cp -a $ROOTFS_DIR/* /mnt/
# 5. Copy kernel:    cp $KERNEL_FILE /mnt/boot/vmlinuz ; cp $INITRAMFS_FILE /mnt/boot/initramfs.img
# 6. Install GRUB:
grub-install --target=i386-pc --boot-directory=/mnt/boot ${LOOPDEV}
cp bootloader/grub.cfg /mnt/boot/grub/grub.cfg
\end{lstlisting}

The \code{--target=i386-pc} selects legacy BIOS GRUB (not UEFI), installed to
the loop device's MBR. The matching \file{bootloader/grub.cfg} is short and
uses the \code{linux}/\code{initrd} commands to load the same two artifacts
the QEMU path used, but now from an on-disk \file{root=/dev/sda1}:

\begin{lstlisting}[language=bash, caption={bootloader/grub.cfg --- the AULinux boot menu entry}]
set timeout=5
set default=0

menuentry "AULinux" {
    linux /boot/vmlinuz root=/dev/sda1 rw quiet init=/sbin/init
    initrd /boot/initramfs.img
}
\end{lstlisting}

\begin{warnbox}[This step is not fully automated and is not idempotent]
\file{install\_grub.sh} requires \code{root}, \code{losetup}, \code{parted},
\code{mkfs.ext4}, and \code{grub-install}. When not run as root it only
\emph{prints} the recipe --- nothing happens to disk. When run as root it
prompts interactively before touching loop devices and mounting under
\file{/mnt}. There is no rollback if a step fails midway, so treat it as an
operator-assisted procedure rather than a hands-off build stage.
\end{warnbox}

\section{Verification (\file{verify.sh})}

\file{verify.sh} is a fast smoke test run against \file{build/} (before any
rootfs is built) that confirms the three core binaries actually execute and
behave. It checks:

\begin{itemize}
  \item \textbf{auutils} --- that \code{pwd} returns the real working
        directory, that \code{echo} round-trips a string, and that \code{ls}
        exits cleanly.
  \item \textbf{aush} --- that the \code{aush} symlink exists, that
        \code{--version} reports ``aush (AULinux Shell)'', and that
        command-mode \code{aush -c "echo ..."} works.
  \item \textbf{aupkg} --- that the binary exists, that \code{--version}
        reports ``aupkg 1.0.0'', and that \code{aupkg list} runs without
        crashing.
\end{itemize}

\begin{lstlisting}[language=bash, caption={verify.sh --- representative behavioral checks}]
# Test command execution mode (-c)
SHELL_EXEC=$("$BUILD_DIR/utils/aush" -c "echo 'Shell Command Execution Mode works!'" 2>&1)
if [ "$SHELL_EXEC" = "Shell Command Execution Mode works!" ]; then
    success "aush command execution (-c) works correctly"
else
    error "aush command execution failed, output: '$SHELL_EXEC'"
    exit 1
fi
\end{lstlisting}

These are genuine functional assertions, not mere ``does the file exist''
checks --- a failed assertion exits non-zero, making \file{verify.sh} suitable
as a CI gate after \code{build.sh all}.

\section{Reproducibility and Limitations}

The pipeline makes a real effort at reproducibility --- \code{-trimpath},
\code{-buildid=}, stripped symbols, \code{CGO\_ENABLED=0},
\code{-extldflags '-static'}, and pinned Rust target triples all reduce
build-to-build variance. But honesty requires naming what is automated versus
manual, and what the scripts assume about the host:

\begin{itemize}
  \item \textbf{Host coupling.} \file{make\_rootfs.sh} hard-codes
        Debian/Ubuntu \code{x86\_64} library paths. Other distributions,
        libc implementations, or architectures will break the library-copy
        stage.
  \item \textbf{No bundled kernel.} AULinux does not build a kernel.
        \file{run\_qemu.sh} and \file{build\_iso.sh} scavenge one from the host
        or download an Alpine \code{vmlinuz-virt}, so a boot may require
        network access and the exact kernel varies by machine.
  \item \textbf{Kernel modules are optional and host-dependent.} They build
        only with matching kernel headers; otherwise the build warns and
        continues without them.
  \item \textbf{Seeded packages are mocks.} The graphics/GUI/Wayland/XFCE/
        desktop-environment/networking repositories contain stub payloads to
        exercise \code{aupkg}, not functional upstream software.
  \item \textbf{GRUB install is operator-assisted.} \file{install\_grub.sh}
        needs root and loop devices and is interactive; the QEMU and ISO paths
        are the fully automated ones.
  \item \textbf{DHCP client is host-sourced.} \file{make\_rootfs.sh} copies
        whatever DHCP client is installed on the build host. If none is found,
        the rootfs boots without a DHCP daemon and network must be configured
        manually.
\end{itemize}

\begin{infobox}[The end-to-end path, in one breath]
\code{run\_qemu.sh} $\rightarrow$ \code{build.sh all} (init, utils/shell,
aupkg, kernel modules into \file{build/}) $\rightarrow$ \code{make\_rootfs.sh}
(FHS tree, static binaries, glibc + host archivers, DHCP client, DNS config,
\code{services.conf} with supervised \code{dhcpcd}, seeded \code{aupkg} repo
into \file{rootfs/}) $\rightarrow$ \code{find | cpio --format=newc | gzip}
(gzipped-cpio initramfs) $\rightarrow$ \code{qemu-system-x86\_64 -kernel
-initrd -append "init=/sbin/init ..."} --- a running, network-capable AULinux
on a serial console, from one command.
\end{infobox}
